
\par\medskip\noindent\textbf{
1 Introduction}\par


Propositional proof complexity investigates proof systems used to demonstrate the unsatisfiability of Boolean formulas.
A central goal – often referred to as Cook’s program, motivated by the $\mathrm{NP}$ vs. $\mathrm{coNP}$ problem –
is to establish superpolynomial lower bounds on the size of refutations within
stronger and stronger proof systems.\par



Resolution is the most
extensively studied such a system, valued for its conceptual simplicity and its
close relationship with modern SAT solvers.
A resolution refutation of a CNF formula $\varphi$ is a sequence of clauses $C_{1},C_{2},\dots,C_{s}$,
concluding with the empty clause (representing a contradiction).
Each clause in the sequence is either an original clause of $\varphi$ or is derived from earlier clauses using
the resolution rule: $\frac{A\lor x\ \ \ B\lor\lnot x}{A\lor B}$.
Numerous techniques have been developed for proving lower bounds in Resolution, and exponential lower bounds are
known for a wide range of formulas. Notably, Urquhart [33] showed that certain unsatisfiable systems of
linear equations over $\mathbb{F}_{2}$ require exponential-size resolution refutations when encoded naturally in CNF.\par



In this paper, we study the proof system resolution over parities ($\mathrm{Res}(\oplus)$), which extends classical resolution
by integrating linear algebra over the finite field $\mathbb{F}_{2}$.
Proof lines in this proof system are disjunctions of linear equations over $\mathbb{F}_{2}$, called
linear clauses.
A $\mathrm{Res}(\oplus)$ refutation of a CNF formula $\varphi$ is a sequence of linear clauses $C_{1},C_{2},\dots,C_{s}$
that ends with the empty clause (a contradiction), every clause of this sequence is either a clause of
$\varphi$ or is obtained from previous clauses by one of the two inference rules:\par


\begin{itemize}


\item 
1. 


The resolution rule, which infers $C\lor D$ from premises $C\lor(f=0)$ and $D\lor(f=1)$
for some linear form $f$.\par





\item 
2. 


The weakening rule, which derives a linear clause $D$ from $C$ if $C$ semantically implies $D$,
i.e. if any assignment satisfying $C$ also satisfies $D$.\par





\end{itemize}



The question of proving a superpolynomial lower bound on the size of $\mathrm{Res}(\oplus)$ refutations remains open
and appears to be very challenging.
The study of lower bounds for $\operatorname{\mathrm{Res}(\oplus)}$ is motivated by the long-standing challenge of proving exponential
lower bounds for Frege systems – a formalization of textbook propositional logic.
Despite decades of effort, no such bounds are known, even for considerably weaker subsystems.
The strongest Frege subsystem for which we currently have lower bounds is $\mathrm{AC}^{0}$-Frege [1],
which operates with constant-depth formulas using only $\land$, $\lor$, and $\lnot$ gates.
However, once parity gates are added – resulting in $\mathrm{AC}^{0}[2]$-Frege – existing lower bound
techniques completely break down.
This is in sharp contrast to circuit complexity, where exponential lower bounds for $\mathrm{AC}^{0}[2]$
circuits have been known for over 30 years [32, 28].
Bridging this discrepancy requires a
deeper understanding of proof systems that reflect the power of reasoning with parities.
As a subsystem of $\mathrm{AC}^{0}[2]$-Frege, $\operatorname{\mathrm{Res}(\oplus)}$ offers a natural and tractable framework for exploring
the power of reasoning with parity, making it a central object of study in this context.\par


\par\medskip\noindent\textbf{
1.1 Recent progress on lower bounds for subsystems of resolution over parities}\par


There are numerous results establishing exponential lower bounds for tree-like $\mathrm{Res}(\oplus)$ refutations of
standard combinatorial formulas, using a variety of techniques [24, 25, 20, 22, 23, 26, 12].\par



Independently, Chattopadhyay, Mande, Sanyal, and Sherif [15] and Beame and Kroth [5] introduced
a lifting approach for establishing lower bounds in tree-like $\mathrm{Res}(\oplus)$.\par



Given a CNF formula $\varphi(y_{1},y_{2},\dots,y_{n})$ and a Boolean function
$g:\{0,1\}^{\ell}\to\{0,1\}$ (called a gadget), we define the
lifted formula $\varphi\circ g$ as the CNF encoding of the formula
$\varphi(g(x_{1,1},x_{1,2},\dots,x_{1,\ell}),\dots,g(x_{n,1},x_{n,2},\dots,x_{n,\ell}))$ where each variable $y_{i}$ in $\varphi$ is replaced by
$g(x_{i,1},\dots,x_{i,\ell})$ for fresh variables $x_{i,1},x_{i,2},\dots,x_{i,\ell}$.\par



Chattopadhyay, Mande, Sanyal, and Sherif [15] introduced the notion of
$k$-stifling gadgets as follows.
A Boolean function $g:\{0,1\}^{\ell}\to\{0,1\}$ is called a $k$-stifling gadget
if, for every $a\in\{0,1\}$ and every choice of $\ell-k$ input variables,
there exists a setting of those $\ell-k$ variables such that the output of $g$
is always equal to $a$, regardless of the values of the remaining $k$ variables.
They further showed that if every resolution refutation of a formula $\varphi$ has
depth at least $h$, and $g$ is a $k$-stifling gadget, then any tree-like $\mathrm{Res}(\oplus)$
refutation of the lifted formula $\varphi\circ g$ must have size at least $2^{kh}$.\par



Efremenko, Garlik, and Itsykson [17] made the first significant progress beyond
the tree-like setting by establishing exponential lower bounds for bottom-regular
$\mathrm{Res}(\oplus)$ refutations of the Binary Pigeonhole Principle formula $\mathrm{BPHP}^{n+1}_{n}$.
Their work introduced the notions of closure and a random-walk technique, both of
which have proven instrumental in subsequent research.
Building on these ideas, and combining them with lifting techniques, Bhattacharya, Chattopadhyay,
and Dvořák [10] showed that certain formulas require exponential-size refutations in
bottom-regular $\mathrm{Res}(\oplus)$, while still admitting polynomial-size refutations in
Resolution.\par



Alekseev and Itsykson [2] showed that one can construct formulas hard for bottom-regular
$\mathrm{Res}(\oplus)$ based on any formula that requires large resolution depth.
Specifically, they proved that if $\varphi$ is an unsatisfiable CNF formula over $n$ variables
with resolution depth at least $\Omega(n)$, then any regular $\mathrm{Res}(\oplus)$ refutation of the
lifted and transformed formula $\mathrm{mix}(\varphi)\circ g$ must have size at least $2^{\Omega(n)}$,
where $g$ is a constant-size gadget and $\mathrm{mix}$ is a semantic-preserving transformation of $\varphi$.\par



Moreover, Alekseev and Itsykson [2] made progress beyond bottom-regular $\mathrm{Res}(\oplus)$ by
establishing a tradeoff between the size and depth of general $\mathrm{Res}(\oplus)$ refutations.
Specifically, they constructed a family of formulas – lifted Tseitin formulas – over $n$ variables and of size $\mathrm{poly}(n)$
such that any $\mathrm{Res}(\oplus)$ refutation must have either size at least $2^{\Omega(n/\log n)}$ or
depth at least $\Omega(n\log\log n)$.
Subsequently, Efremenko and Itsykson [18] improved the depth lower bound to $\Omega(n\log n)$.
In particular, this result implies exponential lower bounds for regular $\mathrm{Res}(\oplus)$ for all
reasonable notions of regularity.
However, a limitation of this result is that it applies to a specific formula.
In this paper, we address this issue by developing a general lifting result.\par



Independently and in parallel with our work, two recent papers appeared that further strengthen depth lower bounds, albeit at the expense of weakening the size bound in the tradeoff described above. Specifically, Bhattacharya and Chattopadhyay [9] proved that for every $\varepsilon>0$, every $\mathrm{Res}(\oplus)$ refutation of lifted Tseitin formulas must either have size $2^{\Omega(n^{\varepsilon})}$ or depth $\Omega(n^{2-\varepsilon})$, where $n$ denotes the number of variables.
Byramji and Impagliazzo [13] established that every Depth-$D$ $\mathrm{Res}(\oplus)$ refutation of the weak binary pigeonhole principle $\mathrm{BPHP}^{m}_{n}$ requires size $2^{\Omega(n^{3}/D^{2})}$. It is worth noting that the depth of any $\mathrm{Res}(\oplus)$ refutation of this formula is at least $\Omega(n)$, so their bound is meaningful for $D$ ranging from $\Omega(n)$ up to $o(n^{3/2})$.\par


\par\medskip\noindent\textbf{
1.2 Main question addressed}\par


It is important to note that, for the lifted Tseitin formulas and the binary pigeonhole principle used in the size-or-depth lower bound of [2, 18, 9] and [13],
it remains unclear whether they actually admit short $\mathrm{Res}(\oplus)$ refutations.
In other words, it is still unclear whether the observed phenomenon constitutes a genuine tradeoff or merely reflects
the current limitations of our techniques for proving size lower bounds.
In this paper, we address the following question: Does there exist a formula that admits a
short $\mathrm{Res}(\oplus)$ refutation, yet any such refutation must have either superlinear
(in the number of variables) depth or exponential (in the size of the formula) size?\par



A negative answer to this question – combined with the result from [2] –
would yield exponential lower bounds on the size of $\operatorname{\mathrm{Res}(\oplus)}$ refutations.\par



On the other hand, a positive answer would establish a supercritical tradeoff
between size and depth in $\operatorname{\mathrm{Res}(\oplus)}$.
Here, supercritical means that for refutations of small size, the required depth
significantly exceeds the worst-case upper bound achievable in the unrestricted setting.
A positive answer would also lend support to a possible explanation for why proving
$\mathrm{Res}(\oplus)$ lower bounds for seemingly simple formulas – such as the pigeonhole
principle – remains so challenging.
It may be that these formultas do admit short refutations, but all such refutations
necessarily have large depth, making them difficult to construct.\par


\par\medskip\noindent\textbf{
1.3 Supercritical tradeoffs in proof complexity}\par


A supercritical tradeoff between two proof complexity measures $\mu$ and $\nu$ for a formula
$\varphi$ occurs when $\varphi$ has proofs with small $\mu$ and others with small $\nu$,
but any proof with $\mu$ below a certain threshold forces $\nu$ to significantly exceed
the worst-case upper bound known for all formulas.\par



In the last few years, many supercritical tradeoffs in proof complexity have been established
[29, 4, 7, 6, 31, 30, 8, 19, 11, 16, 21, 14].
We briefly overview the most relevant results concerning Resolution and $\mathrm{Res}(\oplus)$.
Razborov [29] established a supercritical tradeoff between width and size for tree-like Resolution.
Fleming, Pitassi, and Robere [19] proved supercritical
tradeoffs between size/width and depth for Resolution.
More recently, Buss and Thapen [11] introduced a simple and highly flexible construction
yielding a supercritical tradeoff between size/width and depth in Resolution.
Finally, de Rezende et al. [16] and Göös et al. [21] showed that many of
these tradeoffs can be made truly supercritical, meaning that the lower bounds are expressed in
terms of the formula’s size, rather than merely the number of variables.\par



Chattopadhyay and Dvořák [14] established a supercritical tradeoff between width and
size for tree-like $\mathrm{Res}(\oplus)$.
Their proof builds on a corresponding lifting theorem, which directly lifts Razborov’s result for
tree-like Resolution [29] to the $\mathrm{Res}(\oplus)$ setting.\par



Note that all the tradeoffs mentioned above have been established for proof systems for which
superpolynomial size lower bounds are already known.\par


\par\medskip\noindent\textbf{
1.4 Our contributions}\par


Our main result is the following theorem.\par



\par\medskip\noindent\textbf{
Theorem 1 (Theorem 25).
}\par



Let $\psi$ be an unsatisfiable CNF formula such that $\psi$ does not have
a resolution refutation with width at most $w$ and depth at most $h$.
Assume that there is a natural number $s\geq 2$ such that $h\geq s^{2}w-w$.
Let $g\colon\{0,1\}^{\ell}\to\{0,1\}$ be a 2-stifling gadget (for example,
$g$ can be the 5-bit Majority function).
Then any $\mathrm{Res}(\oplus)$ refutation of $\psi\circ\oplus_{s}\circ g$ has either
size at least $2^{w}$ or depth at least $\frac{s^{2}w}{4\ell}$.\par




The size of the formula $\psi\circ\oplus_{s}\circ g$ is exponential in $s$, so in applications of Theorem 1 we choose $s$ to be as small as possible. On the other hand, to obtain non-trivial depth lower bounds, we need $w$ to be as large as possible.
Two closely related width–depth tradeoff results for resolution are available in the literature, both addressing the regime in which the width is close to the total number of variables $n$: one due to Fleming, Pitassi, and Robere [19], and another due to Buss and Thapen [11].
The former applies for widths $w=n^{1-\epsilon}$, while the latter applies for widths $w=\Omega(n/\log n)$, which is significantly more suitable for our purposes. Applying the lifting from Theorem 1 to the width–depth tradeoff established by Buss and Thapen [11], we obtain a supercritical tradeoff between the size and depth of $\mathrm{Res}(\oplus)$ refutations.\par



\par\medskip\noindent\textbf{
Theorem 2 (Theorem 27).
}\par



For every natural $K\geq 2$ and $n\geq 2$ such that
$K\leq\sqrt{\frac{n}{10\log^{3}n}}$ there is a CNF formula $\varphi_{n,K}$
that contains $O(Kn\log n)$ variables, the formula $\varphi_{n,K}$ is an
$O(K\log^{2}n)$-CNF of size $n^{O(K\log n)}$.
The formula $\varphi_{n,K}$ has a resolution refutation of size $n^{O(K\log n)}$
and of width $O(K\log^{2}n)$ but every $\mathrm{Res}(\oplus)$ refutation of $\varphi_{n,K}$ has either size
at least $2^{\Omega(n/\log n)}$ or depth at least $\Omega(K^{2}n\log n)$.\par




To our knowledge, this is the first instance of a supercritical tradeoff demonstrated
in a proof system lacking known superpolynomial lower bounds on proof size.\par



\par\medskip\noindent\textbf{
Corollary 3 (Corollary 28).
}\par



For every $\delta>0$, Depth-$\frac{n^{4/3}}{\log^{4/3+\delta}n}$
$\mathrm{Res}(\oplus)$ does not p-simulate Resolution.\par




Another important specific case of Theorem 1 is a size-depth
tradeoff for $\mathrm{Res}(\oplus)$ obtained by lifting from resolution width.\par



\par\medskip\noindent\textbf{
Theorem 4 (Theorem 29).
}\par



Let $\psi_{n}$ be a family of unsatisfiable $O(1)$-CNF formulas such that
$\psi_{n}$ has $n$ variables and the resolution width of $\psi_{n}$ is $w(n)$.
For every natural $K\geq 2$ consider a formula
$\Psi_{n,K}:=\psi_{n}\circ\oplus_{K}\circ Maj_{5}$; it has $5nK$ variables,
$\Psi_{n,K}$ is an $O(K)$-CNF formula of size at most $\mathrm{poly}(n)2^{K}$ and any
$\mathrm{Res}(\oplus)$ refutations of $\Psi_{n}$ has either depth at least $\Omega(w(n)K^{2})$
or size at least $2^{\Omega(w(n))}$.\par




Consider several interesting specific cases of Theorem 4.\par



\begin{itemize}


\item 
$\blacksquare$ 


Size-depth tradeoff for formula of size polynomial in number of variables.
Let $w(n)=\Omega(n)$.
Then the formula $\Psi_{n,\lceil\log n\rceil}$ contains $m:=\Theta(n\log n)$
variables, $\Psi_{n,\lceil\log n\rceil}$ is an $O(\log m)$-CNF formula of size
$\mathrm{poly}(m)$ and any $\mathrm{Res}(\oplus)$ refutations of $\Psi_{n,\lceil\log n\rceil}$ has
either depth at least $\Omega(m\log m)$ or size at least $2^{\Omega(m/\log m)}$.\par




This result extends the result from [18] for a large number of formulas.\par





\item 
$\blacksquare$ 


Maximal depth.
Let $w(n)=\Omega(n)$. Consider an arbitrary $1>\delta>0$.
The formula $\Psi_{n,\lceil n^{1-\delta}\rceil}$ contains $m:=\Theta(n^{2-\delta})$
variables, $\Psi_{n,\lceil n^{1-\delta}\rceil}$ is a CNF formula of size
$poly(n)2^{n^{1-\delta}}$ and any $\mathrm{Res}(\oplus)$ refutations of
$\Psi_{n,\lceil n^{1-\delta}\rceil}$ has either depth at least $\Omega(m^{3/2-\delta/2})$
or size at least $2^{\Omega(n)}$.\par




So if we do not restrict ourselves to formulas of polynomial size in the number of variables,
then we can get a superpolynomial size lower bound for depth less than $m^{3/2-\delta}$.\par





\item 
$\blacksquare$ 


Minimal width.
Let $w(n)=\Omega\left(n^{1/2+\delta}\right)$, where $1/2>\delta>0$.
The formula $\Psi_{n,\lceil n^{1/2}\rceil}$ contains $m:=\Theta(n^{3/2})$ variables,
$\Psi_{n,\lceil n^{1/2}\rceil}$ is a CNF formula of size $poly(n)2^{n^{1/2}}$
and any $\mathrm{Res}(\oplus)$ refutations of $\Psi_{n,\lceil n^{1/2}\rceil}$ has either
depth at least $\Omega\left(n^{3/2+\delta}\right)=\Omega\left(m^{1+2\delta/3}\right)$ or size at least
$2^{\Omega\left(n^{1/2+\delta}\right)}$.\par




So we can get a non-trivial size-depth tradeoff starting from a formula with the resolution width
$\Omega\left(n^{1/2+\delta}\right)$.\par





\end{itemize}


\par\medskip\noindent\textbf{
1.5 Technique}\par


Proof of Theorem 1 builds on the size-depth tradeoff established
by Alekseev and Itsykson [2], with subsequent improvements by
Efremenko and Itsykson [18].\par



As a first step, we develop a more flexible size-depth tradeoff that applies to a broad
class of formulas.
Below, we outline the main ideas behind the tradeoff established by Alekseev and
Itsykson [2].\par



Consider a $\mathrm{Res}(\oplus)$ refutation $\Pi$.
We identify certain linear clauses within $\Pi$ as good clauses.
By definition, a good clause cannot be an axiom of the original formula.
These clauses satisfy a crucial property we refer to as the dichotomy property.
Specifically, for every good clause $C$ of moderately small width, one of the following holds:\par


\begin{itemize}


\item 
$\blacksquare$ 


The size of the refutation $\Pi$ is exponential;\par





\item 
$\blacksquare$ 


There exists another good linear clause in $\Pi$ that
appears at a significantly greater depth than $C$ and whose width exceeds that of
$C$ by only a small amount.\par





\end{itemize}



Assuming that the empty clause is good and the size of $\Pi$ is small, the dichotomy property implies that one can iteratively find increasingly deeper good clauses within $\Pi$, ultimately yielding a lower bound on the depth.\par



The dichotomy property is established through a combination of a random walk argument and a bottleneck argument. We begin by defining a set $\Sigma$ of good full assignments that falsify a given linear clause $C_{0}$ (for simplicity, one may think of $\Sigma$ as the set of all assignments falsifying $C_{0}$). We then select a random assignment $\sigma\in\Sigma$ and perform a $t$-step random walk along the refutation graph. At each step, we move from a linear clause to one of its premises that is also falsified by $\sigma$, counting only applications of resolution rules (applications of weakening are ignored).\par



The random walk theorem asserts that, with probability at least $p$, such a walk ends at a good linear clause. Now, if all good linear clauses reachable from $C_{0}$ within $t$ steps have width greater than $|C_{0}|+s$, then there must be at least $p\cdot 2^{s}$ such clauses. This is because a random assignment falsifying $C_{0}$ can falsify a clause of width at least $|C_{0}|+s$ with probability at most $2^{-s}$.\par



Alekseev and Itsykson [2] proved a random walk theorem for formulas of the form $\varphi\circ g$, where $g$ is a 2-stifling gadget and $\varphi$ is an unsatisfiable CNF formula that admits a sufficiently strong strategy for Delayer in a specific Prover–Delayer game. This game is defined with respect to the formula $\varphi$ and a set $\mathcal{A}$ of its partial assignments, which must satisfy two conditions: (1) no assignment in $\mathcal{A}$ falsifies any clause of $\varphi$, and (2) $\mathcal{A}$ is closed under restriction, i.e., any restriction of an assignment in $\mathcal{A}$ also belongs to $\mathcal{A}$.\par



The game proceeds as follows. It begins with some initial assignment $\rho_{0}\in\mathcal{A}$. In each round, Prover selects a variable $x$ and queries its value. Delayer then has two options:\par


\begin{itemize}


\item 
1. 


Pay one black coin to choose a value $a\in\{0,1\}$ and extend the current assignment by setting $x:=a$; or\par





\item 
2. 


Reply with $*$, allowing Prover to choose the value $a$.\par





\end{itemize}


Regardless of the outcome, Delayer earns one white coin for every move. The game terminates as soon as the current assignment no longer belongs to $\mathcal{A}$.\par



The required property of Delayer’s strategy is as follows: for every starting assignment $\rho_{0}\in\mathcal{A}$ in the $(\varphi,\mathcal{A})$-game, there exists a strategy for Delayer that guarantees earning at least $t-|\rho_{0}|$ white coins while spending at most $n$ black coins, where $t$ is significantly larger than $n$. Alekseev and Itsykson [2] provide an example of such a strategy for Tseitin formulas, in which $t$ is, roughly speaking, the number of edges and $n$ is the number of vertices in the underlying graph.\par



Our first observation is that Delayer strategies for such games can be derived from strategies in the Atserias–Dalmau games [3], which characterize resolution width, via a lifting transformation using a parity gate. This simple but powerful idea allows us to establish Theorem 4, thereby completing the first step of our approach.\par



However, formulas with large resolution width are inherently hard for resolution and can therefore only be used to show that small-depth $\mathrm{Res}(\oplus)$ refutations require large size. To establish a supercritical tradeoff, we also need to demonstrate the existence of short refutations with large depth. Our second step is to refine the lifting theorem so that it can be applied starting from formulas whose every resolution refutation must have either width at least $w$ or depth at least $h$. This refinement precisely enables us to apply lifting to the known supercritical tradeoffs between resolution width and depth.\par



In Section 4, we introduce an analogue of the Atserias–Dalmau games tailored to formulas that require resolution width at least $w$ for any resolution proof of depth at most $h$. The properties of winning positions in these games closely resemble those in the original Atserias–Dalmau games, provided we focus on positions within distance $h$ from the empty position. We then apply the parity gadget to these games. Notably, in the proof of the size-versus-depth tradeoff for suitable parameters, it suffices to consider Delayer’s strategy only on positions that remain close to the empty position. This insight enables us to establish the size-depth tradeoff starting from formulas that have no refutations of width at most $h$ and depth at most $w$ simultaneously, thereby proving Theorem 1.\par


\par\medskip\noindent\textbf{
2 Preliminaries}\par

\par\medskip\noindent\textbf{
2.1 Resolution}\par


Let $\varphi$ be an unsatisfiable CNF formula. A resolution refutation of $\varphi$ is a sequence of clauses $C_{1},C_{2},\dots,C_{s}$ such that $C_{s}$ is the empty clause (i.e., identically false) and for every $i\in[s]$ the clause $C_{i}$ is either a clause of $\varphi$ or is obtained from previous clauses by the resolution rule that allows us to derive a clause $C\lor D$ from clauses $C\lor x$ and $D\lor\lnot x$.\par



The size of a resolution refutation is the number of clauses in it.
The depth of a resolution refutation is the length of the longest path between the empty clause and the clause of the original formula.
The width of a resolution refutation is the maximal size of a clause from the refutation.
The resolution width of an unsatisfiable CNF formula $\varphi$ is the minimal possible width
over all resolution refutations of $\varphi$.\par


\par\medskip\noindent\textbf{
2.2 Resolution Over Parities}\par


Here and after, all scalars are from the field $\mathbb{F}_{2}$.
Let $X$ be a set of variables taking values in $\mathbb{F}_{2}$.
A linear form in variables from $X$ is a homogeneous linear polynomial over $\mathbb{F}_{2}$ in variables from $X$ or,
in other words, a polynomial $\sum_{i}^{n}x_{i}a_{i}$, where $x_{i}\in X$ is a variable and $a_{i}\in\mathbb{F}_{2}$ for all $i\in[n]$.
A linear equation is an equality $f=a$, where $f$ is a linear form and $a\in\mathbb{F}_{2}$.\par



A linear clause is a disjunction of linear equations: $\bigvee_{i=1}^{t}(f_{i}=a_{i})$.
Note that over $\mathbb{F}_{2}$ a linear clause $\bigvee_{i=1}^{t}(f_{i}=a_{i})$ may be represented
as the negation of a linear system: $\lnot\bigwedge_{i=1}^{t}(f_{i}=a_{i}+1)$.\par



Now we define the proof system resolution over parities ($\mathrm{Res}(\oplus)$) [25].\par



Let $\varphi$ be an unsatisfiable CNF formula.
A $\mathrm{Res}(\oplus)$ refutation of $\varphi$ is a sequence of linear clauses
$C_{1}$, $C_{2}$, …, $C_{s}$ such that $C_{s}$ is the empty clause (i.e., identically false) and
for every $i\in[s]$ the clause $C_{i}$ is either a clause of $\varphi$ or is obtained from
previous clauses by one of the following inference rules:\par


\begin{itemize}


\item 
$\blacksquare$ 


Resolution rule allows us to derive a linear clause $C\lor D$
from linear clauses $C\lor(f=a)$ and $D\lor(f=a+1)$.\par





\item 
$\blacksquare$ 


Weakening rule allows us to derive from a linear clause $C$ any
linear clause $D$ in the variables of $\varphi$
that semantically follows from $C$ (i.e., any assignment satisfying $C$
also satisfies $D$).\par





\end{itemize}



The size of a $\mathrm{Res}(\oplus)$ refutation is the number of linear clauses in it.
The depth of a $\mathrm{Res}(\oplus)$ refutation is the
maximal number of resolution rules applied on a path between a clause of the initial
formula and the empty clause.\par



\par\medskip\noindent\textbf{
$\blacktriangleright$ Remark 5.
}\par



A resolution refutation of a formula $\varphi$ is a special case of a $\mathrm{Res}(\oplus)$ refutation,
where all linear clauses are plain (i.e., disjunctions of literals).\par




For any function $f(n)$, we denote by Depth-$f(n)$ $\mathrm{Res}(\oplus)$ the subsystem of $\mathrm{Res}(\oplus)$
consisting of refutations with depth at most $f(n)$, where $n$ is the number of variables in
the formula being refuted.\par



For a linear clause $C$ we denote by $L(C)$ the set of linear forms that appear in $C$;
i.e. $L\left(\bigvee_{i=1}^{t}(f_{i}=a_{i})\right)=\{f_{1},f_{2},\dots,f_{t}\}$.
The same notation we use for linear systems: if $\Psi$ is a $\mathbb{F}_{2}$-linear system,
$L(\Psi)$ denotes the set of all linear forms from $\Psi$.\par


\par\medskip\noindent\textbf{
2.3 Lifted formulas}\par


For every CNF formula $\Phi$ over the variables $Y=\{y_{1},y_{2},\dots,y_{m}\}$ and every
Boolean function $g\colon\{0,1\}^{\ell}\to\{0,1\}$ we define a CNF formula
$\Phi\circ g$ with variables
$X=\{x_{i,j}\mid i\in[m],j\in[\ell]\}$ representing in CNF
$\Phi(g(x_{1,1},x_{1,2},\dots,x_{1,\ell}),g(x_{2,1},x_{2,2},\dots,x_{2,\ell}),\dots,g(x_{m,1},x_{m,2},\dots,x_{m,\ell}))$
(i.e. we substitute to every variable of $\Phi$ the function $g$ applied to $\ell$ fresh variables).
Let $\Phi=\bigwedge\limits_{i\in I}C_{i}$, where $C_{i}$ is a clause for all $i\in I$.
For every $i\in[m]$ we denote by $y_{i}\circ g$ the canonical CNF formula representing
$g(x_{i,1},x_{i,2},\dots,x_{i,\ell})$ which has $\ell$ variables in every clause and by
$(\lnot y_{i})\circ g$ the canonical CNF formula representing $\lnot g(x_{i,1},x_{i,2},\dots,x_{i,\ell})$
which has $\ell$ variables in every clause.
Let $C_{i}=l_{i,1}\lor l_{i,2}\lor\dots\lor l_{i,n_{i}}$, where $l_{i}$
is a literal. Then we denote by $C_{i}\circ g$ a CNF formula that represents $l_{i,1}\circ g\lor l_{i,2}\circ g\lor\dots\lor l_{i,n_{i}}\circ g$ as follows: $C_{i}\circ g$ is the conjunction of all clauses of the form $D_{1}\lor D_{2}\lor\dots\lor D_{n_{i}}$, where $D_{j}$ is a clause of $l_{i,j}\circ g$ for all $j\in[m_{i}]$.
And $\Phi\circ g:=\bigwedge\limits_{i\in I}C_{i}\circ g$.\par



We refer to $\Phi\circ g$ as a formula $\Phi$ lifted with a gadget $g$,
to the set $Y=\{y_{1},y_{2},\dots,y_{m}\}$ as a set of unlifted variables,
and to the set $X=\{x_{i,j}\mid i\in[m],j\in[\ell]\}$ as the set of lifted variables.\par



\par\medskip\noindent\textbf{
Lemma 6 (Folklore, see [25], for example).
}\par



Let $g:\{0,1\}^{\ell}\to\{0,1\}$ be a gadget.
If a CNF formula $\varphi$ has a resolution refutation of size $S$ and width $w$,
then the formula $\varphi\circ g$ has a resolution refutation of size
$S2^{O(w\ell)}$ and width $O(w\ell)$.\par



\par\medskip\noindent\textbf{
2.4 Closure and Amortized Closure}\par


For a set of vectors $U$ in a vector space, we denote by $\langle U\rangle$ the linear span of $U$.\par



We consider the set of propositional variables $X=\{x_{i,j}\mid i\in[m],j\in[\ell]\}$.
The variables from $X$ are divided into $m$ blocks by the value of the first index.
The variables $x_{i,1},x_{i,2},\dots,x_{i,\ell}$ form the $i$-th block, for $i\in[m]$.
We may view the set $X$ as the set of lifted variables with respect to a gadget of size $\ell$.\par



Let $F=\{f_{1},f_{2},\dots,f_{n}\}$ be a set of $\mathbb{F}_{2}$-linear forms with variables from $X$.
Consider a coefficient matrix $M$ of $F$: its columns correspond to $X$, and for all $i\in[n]$,
$i$-th row is the coefficient vector of $f_{i}$.
For every $i\in[m]$, let $M_{i}$ consist of matrix columns corresponding to the variables from
the $i$-th block.
Let $I\subseteq[m]$.
We say that $\{M_{i}\}_{i\in I}$ is safe if there are distinct indices $i_{1},i_{2},\dots,i_{t}\in I$
and elements $v_{i_{1}}\in M_{i_{1}},v_{i_{2}}\in M_{i_{2}},\dots,v_{i_{t}}\in M_{i_{t}}$ such that
$v_{i_{1}},v_{i_{2}},\dots,v_{i_{t}}$ is a basis of $\langle\cup_{i\in I}M_{i}\rangle$.\par



A closure [17] of a set of linear forms $F$ is any inclusion-wise minimal
set $S\subseteq[m]$ such that $\{M_{i}\}_{i\in[m]\setminus S}$ is safe.
Informally, the closure indicates the set of the most essential unlifted variables
for the set of linear forms $F$.\par



\par\medskip\noindent\textbf{
Lemma 7 (Uniqueness [17]).
}\par



For any $F$, its closure is unique.\par




We denote the closure of $F$ by $\operatorname{Cl}(F)$.\par



\par\medskip\noindent\textbf{
Lemma 8 ([17]).
}\par



Closure has the following properties.\par


\begin{itemize}


\item 
1. 


If $F\subseteq G$, then $\operatorname{Cl}(F)\subseteq\operatorname{Cl}(G)$;\par





\item 
2. 


$\operatorname{Cl}(F)=\operatorname{Cl}(\langle F\rangle)$;\par





\item 
3. 


$|\operatorname{Cl}(F)|\leq\dim\langle F\rangle$.\par





\end{itemize}




We also require the notion of amortized closure, introduced by Efremenko and Itsykson [18].
Unlike the plain closure, which can grow dramatically with the addition of a single linear form,
the amortized closure is a superset of the plain closure designed to grow more gradually and smoothly.\par



We say that a subset $A\subseteq[m]$ is coverable w.r.t. $\{M_{i}\mid i\in[m]\}$
if for every $i\in A$ there is $v_{i}\in M_{i}$ such the set $\{v_{i}\mid i\in A\}$ is linearly independent.
For subsets $A,B\subseteq[m]$, we say that $A$ is less than $B$ ($A\preceq B$)
if the largest element in the symmetric difference $A\triangle B$ belongs to $B$.\par



An amortized closure of $F$ [18], denoted by $\operatorname{\widetilde{Cl}}(F)$, is the $\preceq$-maximal
subset of $[m]$ that is coverable w.r.t. $\{M_{i}\mid i\in[m]\}$.
It is easy to see that $\operatorname{\widetilde{Cl}}(F)$ does not depend on the permutation of rows in the coefficient
matrix of $F$.\par



\par\medskip\noindent\textbf{
Lemma 9 (Size bound [18]).
}\par



$|\operatorname{\widetilde{Cl}}(F)|\leq\dim\langle F\rangle$.\par



\par\medskip\noindent\textbf{Proof.}\par


$|\operatorname{\widetilde{Cl}}(F)|$ is at most the rank of a coefficient matrix of $F$ that equals $\dim\langle F\rangle$.
$\hfill\blacktriangleleft$\par



\par\medskip\noindent\textbf{
Lemma 10 ([18]).
}\par



$\operatorname{Cl}(F)\subseteq\operatorname{\widetilde{Cl}}(F)$\par




\par\medskip\noindent\textbf{
Lemma 11 (Lipschitz continuity [18]).
}\par



$\operatorname{\widetilde{Cl}}(F)\subseteq\operatorname{\widetilde{Cl}}(F\cup\{f\})$ and $|\operatorname{\widetilde{Cl}}(F\cup\{f\})|\leq|\operatorname{\widetilde{Cl}}(F)|+1$.\par




\par\medskip\noindent\textbf{
Lemma 12 ([18]).
}\par



Let $\Phi$ and $\Psi$ be two linear systems in variables $X=\{x_{i,j}\mid i\in[m],j\in[\ell]\}$.
Let $\pi$ be a partial assignment defined on $\{x_{i,j}\mid i\in\operatorname{Cl}(L(\Phi)),j\in[\ell]\}$.
Let $\Sigma$ consist of all solutions $\sigma$ of $\Phi$ such that $\sigma$ extends $\pi$.
Assume that $\Sigma\neq\emptyset$.
Let $\tau$ be a random element of $\Sigma$.
Then $\Pr[\tau\mbox{ satisfies }\Psi]\leq 2^{|\operatorname{\widetilde{Cl}}(L(\Phi))|-|\operatorname{\widetilde{Cl}}(L(\Psi))|}$.\par



\par\medskip\noindent\textbf{
2.5 Lifting via stifling gadgets}\par


In the lifting settings, we will identify subsets of $[m]$ with corresponding subsets of the
lifted variables $Y$.
It is especially convenient to use such correspondence for closure and amortized closure.
So, we will assume that the support and the (amortized) closure of the set of linear forms over
lifted variables is the set of unlifted variables.\par



A partial assignment $\rho$ to the set of variables $X$ is called block-respectful if,
for every $i$, $\rho$ either assigns values to all variables with support $i$ or does not assign
values to any of them.\par



Suppose that $\rho$ is a block-respectful partial assignment.
Then we define by $\hat{\rho}$ the partial assignment on the set of variables $Y$ such that
$\hat{\rho}(y_{i})=g(\rho(x_{i,1},x_{i,2},\dots,x_{i,\ell}))$ (here we assume that if the
right-hand side is undefined, then the left-hand side is also undefined).\par



Let $k<\ell$. A gadget (i.e. Boolean function) $g:\{0,1\}^{\ell}\to\{0,1\}$ is called
$k$-stifling [15] if for every $A\subset[\ell]$ of size $k$ for every
$c\in\{0,1\}$ there exists $a\in\{0,1\}^{\ell}$ such that for every $b\in\{0,1\}^{\ell}$
if $a$ and $b$ agree on set of indices $[\ell]\setminus A$, then $g(b)=c$.\par



It is easy to see that the majority function $Maj_{2k+1}:\{0,1\}^{2k+1}\to\{0,1\}$
is a $k$-stifling for every $k$.\par



\par\medskip\noindent\textbf{
Lemma 13 ([2]).
}\par



Let $\Psi$ be a satisfiable linear system in the lifted variables $X$.
Let $g\colon\{0,1\}^{\ell}\to\{0,1\}$ be a $1$-stifling gadget.
Suppose there exists a full assignment $\sigma$ to lifted variables $X$ satisfying
$\Psi$ such that $\hat{\sigma}|_{\operatorname{Cl}(L(\Psi))}$ does not falsify any clause of $\varphi$.
Then, $\Psi$ does not contradict any clause of $\varphi\circ g$.\par



\par\medskip\noindent\textbf{
2.6 Supercritical tradeoff between width and depth for resolution}\par


Here, we state the supercritical tradeoff between width and depth in resolution,
as established by Buss and Thapen [11].\par



\par\medskip\noindent\textbf{
Theorem 14 ([11]).
}\par



Let $b,c,d\geq 2$ be integers and $b$ be a power of two.
Then there is an explicit formula $\Phi_{b,c,d}$ that has the following properties:\par


\begin{itemize}


\item 
$\blacksquare$ 


$\Phi_{b,c,d}$ has $bcd+b\log b$ variables;\par





\item 
$\blacksquare$ 


$\Phi_{b,c,d}$ has $c+(d^{c}-1)bc^{2}+b\leq d^{c}b^{2}c$ clauses, of width at most $c+\log b+1$;\par





\item 
$\blacksquare$ 


$\Phi_{b,c,d}$ has a resolution refutation of size $O(d^{c}b^{2}c)$ and width $c+\log b+1$;\par





\item 
$\blacksquare$ 


Any resolution refutation of $\Phi_{b,c,d}$ of width strictly below $b/2$ must have depth at least $d^{c}$.\par





\end{itemize}




\par\medskip\noindent\textbf{
Corollary 15.
}\par



There exists a family of unsatisfiable CNF formulas $\{\Psi_{n}\}_{n=1}^{\infty}$ such that\par


\begin{itemize}


\item 
$\blacksquare$ 


$\Psi_{n}$ contains $n$ variables;\par





\item 
$\blacksquare$ 


the width of $\Psi_{n}$ is $O(\log n)$ and, moreover,
$\Psi_{n}$ has a resolution refutation of size $\mathrm{poly}(n)$ and of
width $O(\log n)$;\par





\item 
$\blacksquare$ 


any resolution refutation of $\Psi_{n}$ of width at most $n/40\log n$
has depth greater than $n^{2}/400\log^{2}n$.\par





\end{itemize}



\par\medskip\noindent\textbf{Proof.}\par


Let $b$ be the maximal power of two such that $5b\log b\leq n$.
Then $n<10b(\log b+1)$. Hence, $b>\frac{n}{20\log n}$.\par



Let us fix $d=2$, $c=2\log b$.\par



Consider $\Phi_{b,c,d}$ from Theorem 14. $\Phi_{b,c,d}$
contains exactly $5b\log b$ variables. We define $\Psi_{n}$ as $\Phi_{b,c,d}$
with $n-5b\log b$ fictive fresh variables.
Precisely we take $\Psi_{n}=\Phi_{b,c,d}\land y_{1}\land y_{2}\dots\land y_{n-5b\log b}$,
where variables $y_{i}$ have no occurrences in $\Phi_{b,c,d}$.\par



It is easy to see that $\Psi_{n}$ contains exactly $n$ variables.
The width of $\Psi_{n}$ is $3\log b+1=O(\log n)$. Any resolution refutation of
$\Phi_{b,c,d}$ is also a refutation of $\Phi_{n}$.
Hence, $\Psi_{n}$ has a resolution refutation of size $\mathrm{poly}(n)$ and of width $O(\log n)$.
Since $\Phi_{b,c,d}$ can be obtained form $\Psi_{n}$ by substitution of all $y_{i}$ to $1$,
we get that any resolution refutation of $\Psi_{n}$ of width at most $n/40\log n$
(which is less than $b/2$) has depth at least $d^{c}=b^{2}>n^{2}/400\log^{2}n$.
$\hfill\blacktriangleleft$\par


\par\medskip\noindent\textbf{
3 Prover-Adversary and Prover-Delayer games}\par


In this section, we define two games based on an unsatisfiable CNF formula $\varphi$.
Let $\mathcal{A}$ be a set of partial assignments for the variables of $\varphi$.
We say that $\mathcal{A}$ is proper for $\varphi$ if the following two properties hold:\par


\begin{itemize}


\item 
$\blacksquare$ 


$\mathcal{A}$ is closed under restrictions: for every $\rho\in\mathcal{A}$
for every $\sigma\subseteq\rho$, $\sigma\in\mathcal{A}$.\par





\item 
$\blacksquare$ 


For every $\sigma\in\mathcal{A}$, $\sigma$ does not falsify any clause of $\varphi$.\par





\end{itemize}


\par\medskip\noindent\textbf{
$(\varphi,\mathcal{A})$-game of Prover and Adversary with starting position $\rho_{0}\in\mathcal{A}$
}\par


In this game, two players, Prover and Adversary, maintain a partial
assignment $\rho$ for variables of $\varphi$ that initially equals $\rho_{0}$.
On every move, Prover chooses a variable $x$, and Adversary earns a coin and chooses a Boolean
value $a$ of $x$.
The current assignment $\rho$ is updated: $\rho:=\rho\cup\{x:=a\}$.
The game ends when $\rho\notin\mathcal{A}$.
The goal of Adversary is to earn as many coins as he can.\par



Let $\mathcal{C}(\varphi)$ denote the set of all partial assignments that do not falsify
any clause of $\varphi$.
It is easy to see that $\mathcal{C}(\varphi)$ is proper and the maximal number of coins that
Adversary can earn in the $(\varphi,\mathcal{C}(\varphi))$ with the empty starting position
is precisely the resolution depth of $\varphi$.
(See [34] for details.)\par



\par\medskip\noindent\textbf{
$\blacktriangleright$ Remark 16.
}\par



In terms of Prover-Adversary games, one can also define the resolution width.
Indeed, Atserias and Dalmau [3] showed that the resolution width of $\varphi$
is at least $w$ if and only if there exists a proper set of assignments $\mathcal{A}$
such that for every $\rho_{0}\in\mathcal{A}$ if $|\rho_{0}|<w$, then in the
$(\varphi,\mathcal{A})$-game with starting position $\rho_{0}$ Adversary has a strategy
that guarantees him to earn at least $|w|-|\rho_{0}|$ coins.\par




Alekseev and Itsykson defined the following games [2].\par


\par\medskip\noindent\textbf{
$(\varphi,\mathcal{A})$-game of Prover and Delayer with starting position $\rho_{0}\in\mathcal{A}$
}\par


In this game, two players, Prover and Delayer, maintain a partial assignment $\rho$ for
variables of $\varphi$ that initially equals $\rho_{0}$.
On every move, Prover chooses a variable $x$, and Delayer has two options:\par


\begin{itemize}


\item 
$\blacksquare$ 


Delayer can earn a white coin and reports $*$.
Then, Prover chooses a Boolean value $a$ of $x$.\par





\item 
$\blacksquare$ 


Delayer can earn a white coin and pay a black
coin to choose a Boolean value $a$ of $x$ by himself.\par





\end{itemize}


The current assignment $\rho$ is updated: $\rho:=\rho\cup\{x:=a\}$. The game ends when $\rho\notin\mathcal{A}$.\par



\par\medskip\noindent\textbf{
$\blacktriangleright$ Remark 17.
}\par



It is easy to see that if, in the $(\varphi,\mathcal{C}(\varphi))$-game starting from
the empty position, there exists a value $t$ such that Delayer has a strategy ensuring
that at some point he accumulates $t$ more white coins than the total number of spent
black coins, then any tree-like resolution refutation of $\varphi$ must have
size at least $2^{t}$. See [27] for details.\par




Delayer’s strategy is called linearly described [2] if there exists a map $f$
that takes as input an ordered set of variables $L$ and a variable $x$, and returns either
$*$ or an $\mathbb{F}_{2}$-affine function $h$ depending on the variables in $L$.
The strategy is applied as follows: given a game history $x_{1}=a_{1},x_{2}=a_{2},\dots,x_{k}=a_{k}$
and a requested variable $x$, Delayer evaluates $f((x_{1},x_{2},\dots,x_{k}),x)$.
If $f((x_{1},x_{2},\dots,x_{k}),x)=*$, then Delayer reports $*$.
Otherwise, if $f((x_{1},x_{2},\dots,x_{k}),x)=h$ for some affine function $h$, Delayer reports
$h(a_{1},a_{2},\dots,a_{k})$.\par


\par\medskip\noindent\textbf{
3.1 Lifting by parity}\par


Let $\mathcal{A}$ be a proper set of partial assignments for a CNF formula
$\varphi(y_{1},y_{2},\dots,y_{n})$.\par



We denote by $\oplus_{k}$ the parity gadget $\{0,1\}^{k}\to\{0,1\}$ that maps
$(a_{1},a_{2},\dots,a_{k})$ to $a_{1}+a_{2}+\dots+a_{k}\bmod 2$.\par



For every partial assignment $\rho$ to the variables of the formula $\varphi\circ\oplus_{k}$
we define the partial assignment $\tilde{\rho}$ to the variables of $\varphi$ as follows:\par


\begin{itemize}


\item 
$\blacksquare$ 


$\tilde{\rho}$ is defined on $y_{i}$, if and only if $\rho$ is defined on all
$x_{i,1},x_{i,2},\dots,x_{i,k}$;\par





\item 
$\blacksquare$ 


$\tilde{\rho}(y_{i})=\bigoplus_{j=1}^{k}\rho(x_{i,j})$.\par





\end{itemize}



Based on the formula $\varphi\circ\oplus_{k}$ we define a set $\mathcal{A}^{\oplus_{k}}$ that
consists of partial assignments $\rho$ to variables of $\varphi\circ\oplus_{k}$ such that
$\tilde{\rho}\in\mathcal{A}$.\par



\par\medskip\noindent\textbf{
Proposition 18.
}\par



If $\mathcal{A}$ is a proper set for $\varphi$, then $\mathcal{A}^{\oplus_{k}}$ is a
proper set for $\varphi\circ\oplus_{k}$.\par



\par\medskip\noindent\textbf{Proof.}\par


Consider $\rho\in\mathcal{A}^{\oplus_{k}}$ and let $\rho^{\prime}\subseteq\rho$.
By the definition, $\tilde{\rho}^{\prime}\subseteq\tilde{\rho}$.
Since $\tilde{\rho}\in\mathcal{A}$,
then $\tilde{\rho}^{\prime}\in\mathcal{A}$, then $\rho^{\prime}\in\mathcal{A}^{\oplus_{k}}$.\par



Consider $\rho\in\mathcal{A}^{\oplus_{k}}$, every clause of $\varphi\circ\oplus_{k}$ is a
clause of $C\circ\oplus_{k}$, where $C$ is a clause of $\varphi$.
Since $\tilde{\rho}$ doesn’t falsify $C$,
there is a variable $y_{j}$ of $C$ such that $\tilde{\rho}$ is not defined on $y_{j}$.
Hence, there is $i\in[k]$ such that $\rho$ is not defined on $x_{j,i}$, hence $\rho$ doesn’t falsify $C\circ\oplus_{k}$.
$\hfill\blacktriangleleft$\par



Now we explore the simple idea of Urquhart [34] that the strategy of
Adversary can be lifted to the strategy of Delayer if we lift the formula by the parity gadget.
The following lemma extends [2, Lemma 6.1].\par



\par\medskip\noindent\textbf{
Lemma 19.
}\par



Assume that Adversary has a strategy in a Prover-Adversary game
$(\varphi,\mathcal{A})$ with starting position $\rho_{0}\in\mathcal{A}$
that guarantees him to earn at least $t$ coins.
Consider a Prover-Delayer game $(\varphi\circ\oplus_{k},\mathcal{A}^{\oplus_{k}})$
with starting position $\sigma_{0}$, where $\tilde{\sigma}_{0}=\rho_{0}$.
Then Delayer has a linearly described strategy that guarantees him to earn at
least $k(t+|\rho_{0}|)-|\sigma_{0}|$ white coins while paying at most $t$ black coins.\par



\par\medskip\noindent\textbf{Proof.}\par


We describe a strategy for Delayer in the Prover–Delayer game $(\varphi\circ\oplus_{k},\mathcal{A}^{\oplus_{k}})$, obtained from the Adversary’s strategy in the Prover–Adversary game
$(\varphi,\mathcal{A})$.
Let $\sigma$ denote the current partial assignment in the first game.
We maintain the invariant that the corresponding partial assignment in the second game is
$\tilde{\sigma}$.\par



Initially, we set $\sigma=\sigma_{0}$ and, thus, $\tilde{\sigma}=\rho_{0}$.
Suppose the requested variable in the first game is $x_{j,i}$.
If there exists an index $i^{\prime}\in[k]\setminus\{i\}$ such that $\sigma$ is undefined on
$x_{j,i^{\prime}}$, then Delayer responds with $*$, and $\tilde{\sigma}$ remains unchanged.\par



Otherwise, if $\sigma$ is defined on all $x_{j,i^{\prime}}$ for $i^{\prime}\in[k]\setminus\{i\}$,
we simulate a Prover request for variable $y_{j}$ in the second game.
Let $a\in\{0,1\}$ be the Adversary’s response according to his strategy.
Delayer then responds with $a\oplus\bigoplus_{i^{\prime}\in[k]\setminus\{i\}}\sigma(x_{j,i^{\prime}})$.\par



To show that this strategy is linearly described, it suffices to prove that the value of $a$
is uniquely determined by the Adversary’s strategy, the ordered list of queried variables,
and the initial assignment $\sigma_{0}$.\par



Indeed, given the ordered list of queried variables in the first game and the initial assignment
$\sigma_{0}$, we can compute both the initial assignment $\rho_{0}$ in the second game and
the corresponding sequence of variable requests.
Since the Adversary’s strategy deterministically specifies the response to each query in
the second game, we can compute all answers, in particular the last one, which is the value of $a$.\par



While $\tilde{\sigma}\in\mathcal{A}$, we have $\sigma\in\mathcal{A}^{\oplus_{k}}$.
Since the Adversary in the first game earns at least $t$ white coins, consider the first
moment when $|\tilde{\sigma}|=|\rho_{0}|+t$.
Each payment of a black coin corresponds to an increment of $\tilde{\sigma}$ by one,
so by this point Delayer has paid exactly $t$ black coins.
The number of earned white coins is at least $|\sigma|-|\sigma_{0}|\geq k|\tilde{\sigma}|-|\sigma_{0}|=k(|\rho_{0}|+t)-|\sigma_{0}|$.
$\hfill\blacktriangleleft$\par


\par\medskip\noindent\textbf{
4 Bounded-depth width games}\par


In this section, we present a combinatorial characterization – an analogue of the
Atserias–Dalmau games [3] – that captures when an unsatisfiable formula
$\varphi$ admits no resolution refutation of both width at most $w$ and depth at most $h$ simultaneously.
A different game characterization for the same property, more closely related to depth-based games,
has already appeared in the literature [7, 19, 18].
In contrast, our characterization is more aligned with width-based games.\par



Let $\varphi$ be a CNF formula. Let $\mathcal{H}$ be a set of pairs $(\rho,i)$ of a
partial assignment $\rho$ and an integer number $i$.
We say that $\mathcal{H}$ is a $(w,h)$-winning strategy for $\varphi$
if the following conditions hold:\par


\begin{itemize}


\item 
$\blacksquare$ 


$(\varepsilon,0)\in\mathcal{H}$, where $\varepsilon$ is an empty assignment.\par





\item 
$\blacksquare$ 


If $(\rho,i)\in\mathcal{H}$, then $|\rho|\leq w$, $i\leq h$ and $\rho$
doesn’t falsify any clause of $\varphi$.\par





\item 
$\blacksquare$ 


If $(\rho,i)\in\mathcal{H}$ and $\rho^{\prime}\subseteq\rho$, then
$(\rho^{\prime},i)\in\mathcal{H}$.\par





\item 
$\blacksquare$ 


If $(\rho,i)\in\mathcal{H}$, $|\rho|<w$, $i<h$,
and $x\in\mathrm{Vars}(\varphi)\setminus\mathrm{Dom}(\rho)$, then there exists $a\in\{0,1\}$
such that $(\rho\cup\{x:=a\},i+1)\in\mathcal{H}$.\par





\end{itemize}



\par\medskip\noindent\textbf{
Theorem 20.
}\par



Let $w\geq 0$ and $h\geq 0$ be some integers;
and let $\varphi$ be an unsatisfiable CNF formula such that $\varphi$ doesn’t have a resolution refutation of width
at most $w$ and simultaneously with depth at most $h$.
Then there exists a $(w,h)$-winning strategy for $\varphi$.\par



\par\medskip\noindent\textbf{Proof.}\par


Proof by induction on $h$. The base case is $h=0$. In this case, we can take $\mathcal{H}$ consisting of the only
element $(\varepsilon,0)$, where $\varepsilon$ is an empty clause.\par



We only have to verify that the formula does not contain an empty clause; this is true since the formula $\varphi$
does not have refutation with depth zero and width zero.\par



Induction step. Let $\phi^{\prime}$ be a CNF formula containing all clauses that can be derived from $\phi$ in at most one step with width at most $w$.
It is easy to see that any resolution refutation of $\phi^{\prime}$ has either width greater than $w$ or depth greater than $h-1$.
We apply the induction hypothesis to $\phi^{\prime}$. Let $\mathcal{H^{\prime}}$ be a $(w,h-1)$ winning strategy for $\phi^{\prime}$.\par


\begin{itemize}


\item 
$\blacksquare$ 


For every $(\rho,i)\in\mathcal{H^{\prime}}$, $\rho$ does not falsify clauses of $\varphi^{\prime}$ and hence clauses of $\varphi$
of width at most $w$; $|\rho|\leq w$, hence $\rho$ does not falsify any clause of $\varphi$.\par





\item 
$\blacksquare$ 


Consider $(\rho,h-1)\in\mathcal{H^{\prime}}$ such that $|\rho|\leq w-1$. Let $x$ be a variable from $\mathrm{Vars}(\varphi)\setminus\mathrm{Dom}(\rho)$.
We claim that either $\rho\cup\{x:=0\}$ or $\rho\cup\{x:=1\}$ does not falsify any clause of $\varphi$.
Suppose that $\rho\cup\{x:=0\}$ falsifies $C_{0}$ and $\rho\cup\{x:=0\}$ falsifies $C_{1}$, where $C_{0}$ and $C_{1}$ are clauses of $\varphi$.
Since $|\rho\cup\{x:=0\}|\leq w$, width of $C_{0}$ and $C_{1}$ are at most $w$.
Since $\rho$ doesn’t falsify neither $C_{0}$ nor $C_{1}$,
$C_{0}=D_{0}\lor x$ and $C_{1}=D_{1}\lor\lnot x$ and $\rho$ falsifies $D_{0}$ and $D_{1}$.
Therefore $\rho$ falsifies $D_{0}\lor D_{1}$.
Note that $|D_{0}\lor D_{1}|\leq|\rho|\leq w$ and $D_{0}\lor D_{1}$ is the result of resolution rule applied to $C_{0}$ and $C_{1}$,
hence, $D_{0}\lor D_{1}$ is a clause of $\phi^{\prime}$.
We get a contradiction since $(\rho,h-1)\in\mathcal{H^{\prime}}$.\par





\end{itemize}


Let us define $\Gamma$ to be the set of all pairs of the form $(\rho\cup\{x:=a\},h)$ such that
$(\rho,h-1)\in\mathcal{H^{\prime}}$, $|\rho|\leq w-1$, and $a\in\{0,1\}$, provided that $\rho\cup\{x:=a\}$
does not falsify any clause of $\varphi$.
Let $\Gamma^{\prime}$ be the set of all pairs $(\rho^{\prime},h)$ such that there exists an assignment
$\rho$ with $\rho^{\prime}\subseteq\rho$ and $(\rho,h)\in\Gamma$.\par



We define $\mathcal{H}:=\mathcal{H^{\prime}}\cup\Gamma^{\prime}$.
Let us verify that $\mathcal{H}$ is a $(w,h)$-winning strategy for $\varphi$.\par


\begin{itemize}


\item 
$\blacksquare$ 


Consider a pair $(\rho,i)\in\mathcal{H}$.
If $i<h-1$, then $(\rho,i)\in\mathcal{H^{\prime}}$, so $|\rho|\leq w$ and $\rho$ does not falsify any clause of $\varphi$.
If $i=h$, then $(\rho,h)\in\Gamma^{\prime}$, which implies that there exists an assignment $\sigma$ such that
$\rho\subseteq\sigma$ and $(\sigma,h)\in\Gamma$.
Therefore, $\sigma$ (and thus $\rho$) does not falsify any clause of $\varphi$, and $|\rho|\leq|\sigma|\leq w$.\par





\item 
$\blacksquare$ 


Consider a pair $(\rho,i)\in\mathcal{H}$ and let $\rho^{\prime}\subseteq\rho$.
If $i<h-1$, then $(\rho,i)\in\mathcal{H^{\prime}}$ and thus $(\rho^{\prime},i)\in\mathcal{H^{\prime}}\subseteq\mathcal{H}$.
If $i=h$, $(\rho,i)\in\Gamma^{\prime}$, hence $(\rho^{\prime},i)\in\Gamma^{\prime}\subseteq\mathcal{H}$.\par





\item 
$\blacksquare$ 


Consider a pair $(\rho,i)\in\mathcal{H}$ such that $|\rho|<w$ and $i<h$, and let $x\in\mathrm{Vars}(\varphi)\setminus\mathrm{Dom}(\rho)$.
If $i<h-1$, then $(\rho,i)\in\mathcal{H^{\prime}}$.
By the properties of $\mathcal{H^{\prime}}$, there exists $a\in{0,1}$ such that $\rho\cup\{x:=a\}$ does not falsify any clause
of $\varphi^{\prime}$ (and hence none of $\varphi$), and $(\rho\cup\{x:=a\},i+1)\in\mathcal{H^{\prime}}\subseteq\mathcal{H}$.
If $i=h-1$, then as noted above, there exists $a\in{0,1}$ such that $\rho\cup\{x:=a\}$ does not falsify any clause
of $\varphi$.
Thus, $(\rho\cup\{x:=a\},i+1)\in\Gamma\subseteq\Gamma^{\prime}\subseteq\mathcal{H}$.\par





\end{itemize}


$\hfill\blacktriangleleft$\par



\par\medskip\noindent\textbf{
Proposition 21.
}\par



Let $\mathcal{H}$ be a $(w,h)$-winning strategy for the formula $\varphi$.
Consider a set $\mathcal{B}$ that consists of all partial assignments $\rho$ such that $(\rho,i)\in\mathcal{H}$ for some $i$.
Then for every $(\rho_{0},i_{0})\in\mathcal{H}$ in the Prover-Adversary game $(\varphi,\mathcal{B})$ there is a strategy of Adversary with starting position $\rho_{0}$ that guarantees him to earn $\min\{w-|\rho_{0}|,h-i_{0}\}$ coins.\par



\par\medskip\noindent\textbf{Proof.}\par


Let $\rho$ denote the current partial assignment in the Prover-Adversary game
$(\varphi,\mathcal{B})$.
Adversary will maintain the number $i$ such that $(\rho,i)\in\mathcal{A}$.
Initially $\rho:=\rho_{0}$ and $i:=i_{0}$.
Let $|\rho|<w$ and $i<h$ and $x$ be the requested variable.
Then there exists such $a\in\{0,1\}$ such that $(\rho\cup\{x:=a\},i+1)\in\mathcal{H}$,
hence $\rho\cup\{x:=a\}\in\mathcal{B}$.
The Adversary chooses any of such $a$ and updates $\rho:=\rho\cup\{x:=a\}$ and $i:=i+1$.
After each step, the value $\min\{w-|\rho|,h-i\}$ decreases by one.
And we can not make the next step if $\min\{w-|\rho|,h-i\}=0$.
Thus Adversary earns $\min\{w-|\rho_{0}|,h-i_{0}\}$ coins.
$\hfill\blacktriangleleft$\par



\par\medskip\noindent\textbf{
Definition 22.
}\par



Consider a $(\varphi,\mathcal{A})$-game between Prover and Delayer, and let $H$ be a strategy for Delayer in this game.
For any two assignments $\sigma,\sigma^{\prime}\in\mathcal{A}$, we define the distance between them, denoted $\Delta_{\mathcal{A},H}(\sigma,\sigma^{\prime})$, as the minimal integer $K$ such that there exists a sequence of assignments $\sigma_{0}=\sigma,\sigma_{1},\dots,\sigma_{n}=\sigma^{\prime}\in\mathcal{A}$ satisfying the following conditions:\par


\begin{itemize}


\item 
$\blacksquare$ 


For every $i\in[n]$, either $\sigma_{i+1}$ is obtained from $\sigma_{i}$ by one step of Delayer’s strategy $H$, or $\sigma_{i+1}\subseteq\sigma_{i}$;\par





\item 
$\blacksquare$ 


The total number of steps of the first type (i.e., applications of $H$) is exactly $K$.\par





\end{itemize}


If no such sequence exists, we define $\Delta_{\mathcal{A},H}(\sigma,\sigma^{\prime})=\infty$.\par




\par\medskip\noindent\textbf{
Proposition 23.
}\par



Let $\mathcal{H}$ be a $(w,h)$-winning strategy for a formula $\varphi$.
Define $\mathcal{B}$ as the set of all partial assignments $\rho$ such that $(\rho,i)\in\mathcal{H}$ for some $i$.
Now consider a strategy $H$ for Delayer in the game $(\varphi\circ\oplus_{k},\mathcal{B}^{\oplus_{k}})$, obtained via Lemma 19 from the Adversary’s strategy in the $(\varphi,\mathcal{B})$-game constructed in Proposition 21.\par


\begin{itemize}


\item 
1. 


Let $(\rho_{0},j)\in\mathcal{H}$, consider $\sigma_{0}\in\mathcal{B}^{\oplus_{k}}$ such that
$\tilde{\sigma}_{0}=\rho_{0}$. Let $\sigma\in\mathcal{B}^{\oplus_{k}}$ such that
$\Delta_{\mathcal{B}^{\oplus_{k}},H}(\sigma_{0},\sigma)<\infty$.
Then $(\tilde{\sigma},j^{\prime})\in\mathcal{H}$ for some
$j^{\prime}\leq j+\Delta_{\mathcal{B}^{\oplus_{k}},H}(\sigma_{0},\sigma)$.\par





\item 
2. 


Let $\sigma\in\mathcal{B}^{\oplus_{k}}$ such that $\Delta_{\mathcal{B}^{\oplus_{k}},H}(\epsilon,\sigma)\leq h-w$, where $\epsilon$ is an empty assignment. Then, the strategy $H$ with starting position $\sigma$ guarantees that Delayer earns at least $kw-|\sigma|$ white coins while paying at most $w$ black coins.\par





\end{itemize}



\par\medskip\noindent\textbf{Proof.}\par


\begin{itemize}


\item 
1. 


Let us denote $K=\Delta_{\mathcal{B}^{\oplus_{k}},H}(\sigma_{0},\sigma)$.
There are exist $\sigma_{1},\sigma_{2},\dots,\sigma_{m}=\sigma$ such that
for every $i\in[m]$ either
$\sigma_{i}\subseteq\sigma_{i-1}$ or $\sigma_{i}$ can be obtained from $\sigma_{i-1}$ by one step according the strategy $H$.
Let us define the sequence $j_{0},j_{1},\dots,j_{m}$ as follows: $j_{0}=j$, for every $i\in[m]$:
if $\sigma_{i}\subseteq\sigma_{i-1}$, then $j_{i}=j_{i-1}$;
if $\sigma_{i}$ is obtained from $\sigma_{i-1}$ by one step according the strategy $H$, then if $\tilde{\sigma}_{i}=\tilde{\sigma}_{i-1}$,
then $j_{i}=j_{i-1}$ and if $\tilde{\sigma}_{i}\neq\tilde{\sigma}_{i-1}$, then $j_{i}=j_{i-1}+1$.
Taking into account the definition of the game $(\varphi\circ\oplus_{k},\mathcal{B}^{\oplus_{k}})$ and the construction of the strategy $H$, it is easy to verify that $(\tilde{\sigma}_{i},j_{i})\in\mathcal{H}$ for all $i\in[m]$. In particular, this implies that $(\tilde{\sigma},j_{m})\in\mathcal{H}$.
Observe that $j_{i}>j_{i-1}$ only when $s_{i}$ is obtained from $s_{i-1}$ by a single step of the strategy $H$. Therefore, $j_{m}\leq j+K$.\par





\item 
2. 


By the previous item, there exists $j^{\prime}\leq h-w$ such that $(\tilde{\sigma},j^{\prime})\in\mathcal{H}$. Proposition 21 gives a strategy of Adversary in the Prover-Adversary game $(\varphi,\mathcal{B})$ with starting position $\tilde{\sigma}$ that guarantees him to earn $w-|\tilde{\sigma}|$ coins. Then by Lemma 19 the strategy $H$ in the Prover-Delayer game with starting position $\sigma$ guarantees Delayer to earn at least $k(w-|\tilde{\sigma}|+|\tilde{\sigma}|)-|\sigma|=kw-|\sigma|$ white coins while paying at most $w$ black coins.\par





\end{itemize}


$\hfill\blacktriangleleft$\par


\par\medskip\noindent\textbf{
5 Random-walk theorem}\par


In this section, we present the main tool developed by Itsykson and Alekseev [2] for proving size-depth tradeoffs in $\mathrm{Res}(\oplus)$.\par



Let $\Pi$ be a $\mathrm{Res}(\oplus)$ refutation, $C_{0}$ be a linear clause from $\Pi$, $\Sigma$ be a set of full assignments that falsify $C_{0}$, and
$t\in\mathbb{N}$ be a natural number.
A $(\Pi,C_{0},\Sigma,t)$-random walk is defined as follows: sample an assignment $\sigma$ uniformly at random from $\Sigma$,
and perform a walk of weighted length $t$ on the refutation graph of $\Pi$, starting at the node labeled by $C_{0}$.
At each step, the walk proceeds from a linear clause to a premise falsified by $\sigma$.
There are two cases: if the step uses the weakening rule, there is a single premise and the step has weight zero;
if it uses the resolution rule, there are two premises and the step has weight one.
The walk terminates either upon reaching a clause from the initial formula or when the total weight accumulated over all steps reaches $t$.
If the walk terminates at a node labeled with a linear clause $C$, then $C$ is the value of the random variable defined by the walk.\par



\par\medskip\noindent\textbf{
Theorem 24 (Theorem 4.3 from [2]).
}\par



Let $\varphi$ be an unsatisfiable CNF formula and $g\colon\{0,1\}^{\ell}\to\{0,1\}$ be a 2-stifling gadget.
Consider a $\mathrm{Res}(\oplus)$ refutation $\Pi$ of $\varphi\circ g$ and a linear clause $C_{0}$ from $\Pi$.
Let $\tau$ be a solution of $\lnot C_{0}$ and let $\rho_{0}$ be the restriction of $\hat{\tau}$ to $\operatorname{Cl}(L(C_{0}))$.
Let $\Sigma$ be the set of all full assigmnets $\pi$ such that $\pi$ satisfies $\lnot C_{0}$ and $\hat{\pi}$ extends $\rho_{0}$.
Let $t$ be integer number such that $t\leq w-|\operatorname{\widetilde{Cl}}(L(C_{0}))|+|\rho_{0}|$. Let a linear clause $C$ denote the result of the $(\Pi,C_{0},\Sigma,t)$-random walk defined by a random variable $\sigma$ distributed uniformly on $\Sigma$.\par




Let $\mathcal{A}$ be a proper set of partial assignments for $\mathrm{Vars}(\varphi)$.
Assume that in the $(\varphi,\mathcal{A})$-game with starting position $\rho_{0}\in\mathcal{A}$,
Delayer has a linearly described strategy $H$ that guarantees him to earn $w$ white coins while paying at most $c$ black coins.
Then $\hat{\sigma}|_{\operatorname{Cl}(L(C))}\in\mathcal{A}$ and
$\Delta_{\mathcal{A},H}(\rho_{0},\hat{\sigma}|_{\operatorname{Cl}(L(C))})\leq w$ with probability at least $2^{-c(\ell-1)}$.\par




Theorem 24 is a slightly modified version of [2, Theorem 4.3],
with two minor adjustments to the statement that do not affect the validity of the original proof.
The first modification, introduced in [18], concerns the inequality for $t$, namely the bound
$t\leq w-|\operatorname{\widetilde{Cl}}(L(C_{0}))|+|\rho_{0}|$ was originally stated in a stronger form as
$t\leq w-\mathrm{rk}{(\lnot C_{0})}+|\rho_{0}|$.
We refer the reader to [18] for an explanation of why the proof remains valid under the weaker bound.\par



Here, we focus on the second modification: namely, we additionally assert that
$\Delta_{\mathcal{A},H}(\rho_{0},\hat{\sigma}|_{\operatorname{Cl}(L(C))})\leq w$.\par



First, observe that in the “lucky” execution of the random walk – that is, when
$\hat{\sigma}|_{\operatorname{Cl}(L(C))}\in\mathcal{A}$ – Lemma 13 implies that $C$ is not a clause of $\varphi\circ g$.
Hence, the random walk successfully makes $t$ steps (i.e., it does not terminate prematurely at a leaf).
Let us denote the sequence of visited linear clauses by $C_{0},C_{1},\dots,C_{m}=C$.
There exist indices $0\leq i_{1}<i_{2}<\dots<i_{t}\leq m$ such that:\par


\begin{itemize}


\item 
$\blacksquare$ 


For every $i\in\{i_{1},i_{2},\dots,i_{t}\}$, the clause $C_{i}$ is obtained from $C_{i+1}$ and
another premise by applying the resolution rule over the linear form $f_{i}$.\par





\item 
$\blacksquare$ 


For every $i\in\{0,1,\dots,m-1\}\setminus\{i_{1},i_{2},\dots,i_{t}\}$,
the clause $C_{i}$ is obtained from $C_{i+1}$ by the weakening rule.\par





\end{itemize}


By the definitions of the resolution and weakening rules, it follows that
$L(C)\subseteq\langle L(C_{0}),f_{1},f_{2},\dots,f_{t}\rangle$. Hence, by Lemma 8, we have\par





 
\[\begin{gathered}\operatorname{Cl}(L(C))\subseteq\operatorname{Cl}\left(L(C_{0})\cup\{f_{1},f_{2},\dots,f_{t}\}\right).\end{gathered}\] 
 




The proof of [2, Theorem 4.3] shows that, with probability at least $2^{-c(\ell-1)}$, the restriction of
$\hat{\sigma}$ to $\operatorname{Cl}(L(C_{0})\cup\{f_{1},f_{2},\dots,f_{t}\})$ is consistent with the strategy $H$ in the game
$(\varphi,\mathcal{A})$ starting from position $\rho_{0}$.
This means that there exists some sequence of moves by Prover for which Delayer’s strategy $H$ reaches the assignment
$\hat{\sigma}|_{\operatorname{Cl}(L(C_{0})\cup\{f_{1},f_{2},\dots,f_{t}\})}$.
Since $\mathcal{A}$ is closed under restrictions, it follows that $\hat{\sigma}|_{\operatorname{Cl}(L(C))}\in\mathcal{A}$.
The number of moves in the game does not exceed\par





 
$\displaystyle|\operatorname{Cl}\left(L(C_{0})\cup\{f_{1},f_{2},\dots,f_{t}\}\right)|$ 
$\displaystyle\mathop{\leq}^{\mbox{\scriptsize(Lem.\nobreak\ 10)}}|\operatorname{\widetilde{Cl}}\left(L(C_{0})\cup\{f_{1},f_{2},\dots,f_{t}\}\right)|$ 
 




 
 
$\displaystyle\mathop{\leq}^{\mbox{\scriptsize(Lem.\nobreak\ 11)}}|\operatorname{\widetilde{Cl}}(L(C_{0}))|+t\leq w.$ 
 




Hence, $\Delta_{\mathcal{A},H}(\rho_{0},\hat{\sigma}|_{\operatorname{Cl}(L(C))})\leq w$ by the definition of the distance $\Delta$.\par


\par\medskip\noindent\textbf{
6 Size vs depth tradeoff}\par


\par\medskip\noindent\textbf{
Theorem 25.
}\par



Let $\psi$ be an unsatisfiable CNF formula such that $\psi$ does not have a resolution
refutation with width at most $w$ and simultaneously depth at most $h$. Assume that there is a natural number $s\geq 2$ such that $h\geq s^{2}w-w$.
Let $\varphi:=\psi\circ\oplus_{s}$.\par




Let $g\colon\{0,1\}^{\ell}\to\{0,1\}$ be a 2-stifling gadget.
Then any $\mathrm{Res}(\oplus)$ refutation of $\varphi\circ g$ has either size at least $2^{w}$ or depth at least
$\frac{s^{2}w}{4\ell}$.\par



\par\medskip\noindent\textbf{Proof.}\par


By Theorem 20 for the formula $\psi$ there is a $(w,h)$-winning strategy $\mathcal{H}$.
Let $\mathcal{B}$ be the set of all partial assignments $\rho$ such that $(\rho,i)\in\mathcal{H}$ for some $i$.
By Proposition 21, for every $(\rho,i)\in\mathcal{H}$ in Prover-Adversary game
$(\psi,\mathcal{B})$ there is a strategy of Adversary with starting position $\rho$ that guarantees him to earn
$\min\{w-|\rho|,h-i\}$ coins.\par



We define $\mathcal{A}=\mathcal{B}^{\oplus_{k}}$ and consider a linearly described strategy $H$ for Delayer in the
Prover-Delayer game $(\varphi,\mathcal{A})$ that exists by Lemma 19.\par



Let us denote $t:=ws$.
By Proposition 23, for every $\rho_{0}\in\mathcal{A}$ such that $\Delta_{\mathcal{A},H}(\epsilon,\sigma)\leq h-w$,
in the game with starting position $\rho_{0}$ the strategy $H$ guaranties Delayer to earn at least $t-|\rho_{0}|$ white coins
while paying at most $w$ black coins.\par



In what follows, we use our lifting notations assuming that variables of $\varphi$ are unlifted and variables of $\varphi\circ g$ are lifted.\par



Let $C$ be a linear clause over lifted variables (i.e., variables of the formula $\varphi\circ g$),
and let $\rho\in\mathcal{A}$ be a partial assignment over the original (unlifted) variables.
We say that $C$ corresponds to $\rho$ if there exists an assignment $\tau$ satisfying $\lnot C$ such that
the restriction of $\hat{\tau}$ onto $\operatorname{Cl}(L(C))$ coincides with $\rho$, that is, $\hat{\tau}|_{\operatorname{Cl}(L(C))}=\rho$.
We denote this relation by $C\sim\rho$.
For convenience, we also define a measure $\mu$ of a clause $C$ as
$\mu(C):=\left|\operatorname{\widetilde{Cl}}(L(C))\right|$.\par



Consider a $\mathrm{Res}(\oplus)$ refutation of $\varphi\circ g$ and denote it by $\Pi$.\par



\par\medskip\noindent\textbf{
Claim 26. 
}\par



Assume that $\Pi$ contains a linear clause $C_{0}$ such that\par


\begin{itemize}


\item 
$\blacksquare$ 


$\mu(C_{0})\leq r$, where $r<t$;\par





\item 
$\blacksquare$ 


there exists $\rho_{0}\in\mathcal{A}$ such that $C_{0}\sim\rho_{0}$ and $\Delta_{\mathcal{A},H}(\epsilon,\rho_{0})\leq h-w$.\par





\end{itemize}


Let $S_{t-r}(C_{0})$ denote the set of all linear clauses $C$ from $\Pi$ such that\par


\begin{itemize}


\item 
$\blacksquare$ 


there is a path from $C_{0}$ to $C$ of weighted length $t-r$ in the graph of $\Pi$
(computing length, we compute weakening rules with weight zero and resolution rules with weight one);\par





\item 
$\blacksquare$ 


there exists $\rho\in\mathcal{A}$ such that $C\sim\rho$ and $\Delta_{\mathcal{A},H}(\rho_{0},\rho)\leq t-r$.\par





\end{itemize}




Assume that for every $C\in S_{t-r}(C_{0})$, $\mu(C)\geq r+w\ell$. Then, the size of the refutation $\Pi$ is at least $2^{w}$.\par



\par\medskip\noindent\textbf{
Proof.}\par


Notice that $|\rho_{0}|=|\operatorname{Cl}(L(C_{0}))|\mathop{\leq}\limits^{\mbox{\scriptsize(Lemma~10)}}|\operatorname{\widetilde{Cl}}(L(C_{0}))|=\mu(C_{0})\leq r$.\par



Let $\Sigma$ be the set of all assignments $\pi$ such that $\pi$ satisfies $\lnot C_{0}$ and $\hat{\pi}|_{\operatorname{Cl}(L(C_{0}))}=\rho_{0}$. Since $C_{0}\sim\rho_{0}$, $\Sigma\neq\emptyset$.\par



Let a linear clause $C$ denote the result of the $(\Pi,C_{0},\Sigma,t-r)$-random walk defined by a random variable $\sigma$
distributed uniformly on $\Sigma$.
Notice that $t-r\leq(t-|\rho_{0}|)+|\rho_{0}|-|\operatorname{\widetilde{Cl}}(L(C_{0}))|$. Let $\rho=\hat{\sigma}|_{\operatorname{Cl}(L(C))}$.
By Theorem 24, with probability at least $2^{-(\ell-1)w}$, $\rho\in\mathcal{A}$,
$C\sim\rho$ and $\Delta_{\mathcal{A},H}(\rho_{0},\rho)\leq(t-r)$.
By Lemma 13, $C$ is not a clause of $\phi\circ g$,
hence, the length of the path between $C_{0}$ and $C$ is exactly $t-r$, hence $C\in S_{t-r}(C_{0})$.
Thus, $\mu(C)\geq r+w\ell$.\par



Consider some linear clause $D$ such that $\mu(D)\geq r+w\ell$.
By Lemma 12,
$\Pr_{\sigma\in\Sigma}[\mbox{$\sigma$ satisfies $\lnot D$}]\leq 2^{|\operatorname{\widetilde{Cl}}(L(C_{0}))|-|\operatorname{\widetilde{Cl}}(L(D))|}=2^{\mu(C_{0})-\mu(D)}\leq 2^{-w\ell}$.\par



Hence, the refutation $\Pi$ contains at least $\frac{2^{-(\ell-1)w}}{2^{-\ell w}}=2^{w}$ clauses $D$ with $\mu(D)\geq r+w\ell$.
$\hfill\vartriangleleft$\par



Assume that the size of $\Pi$ is less than $2^{w}$. Our goal is to show that under this assumption, the depth of
$\Pi$ is at least $\frac{ws^{2}}{4\ell}$.\par



Let $D_{0}$ denote the empty clause from $\Pi$.
$D_{0}\sim\rho_{0}$, where $\rho_{0}$ equals the empty assignment $\epsilon$ and, thus, $\rho_{0}\in\mathcal{A}$.
Since $|\Pi|<2^{w}$, by Claim 26, there is a clause $D_{1}$ in $\Pi$ such that there is a path from
$D_{0}$ to $D_{1}$ of length $t$ and $\mu(D_{1})\leq w\ell$ and there is $\rho_{1}\in\mathcal{A}$ such that
$D_{1}\sim\rho_{1}$ and $\Delta_{\mathcal{A},H}(\epsilon,\rho_{1})\leq t$.\par



Let $k:=\lfloor\frac{t}{2w\ell}\rfloor$, then $w\ell k\leq t/2$.
We repeat the same reasoning $k-1$ more times for all $i$ from $1$ to $k-1$, maintaining invariant $\mu(D_{i})|\leq w\ell i$.
Since $|\Pi|<2^{w}$, by Claim 26 there is a linear clause $D_{i+1}$ in $\Pi$ such that there is a path from $D_{i}$
to $D_{i+1}$ of length $t-w\ell i$ and there is $\rho_{i}\in\mathcal{A}$ such that $D_{i}\sim\rho_{i}$ and
$\Delta_{\mathcal{A},H}(\rho_{i},\rho_{i+1})\leq t$ and $\mu(D_{i+1})|\leq w\ell(i+1)$.
Note that for all $i\in[k-1]$, by triangle inequalities,
$\Delta(\epsilon,\rho_{i})\leq kt\leq\frac{t^{2}}{2w\ell}=\frac{ws^{2}}{2\ell}\leq h-w$;
the last inequality holds since by the conditions of the theorem $s\geq 2$ and
$h\geq w(s^{2}-1)\geq w(\frac{s^{2}}{2\ell}+1)=\frac{ws^{2}}{2\ell}+w$. So the distance conditions in applications of
Claim 26 are satisfied.\par



So under the assumption $|\Pi|<2^{w}$ we get that the depth of $\Pi$ is at least the length of the path from
$D_{0}$ to $D_{1}$, from $D_{1}$ to $D_{2}$, etc, from $D_{k-1}$ to $D_{k}$ which is at least $kt/2\geq\frac{ws^{2}}{4\ell}$.
$\hfill\blacktriangleleft$\par



We now examine two specific cases of Theorem 25.\par



\par\medskip\noindent\textbf{
Theorem 27.
}\par



For every natural $K\geq 2$ and $n\geq 2$ such that $K\leq\sqrt{\frac{n}{10\log^{3}n}}$ there is a CNF formula
$\varphi_{n,K}$ that contains $5Kn\lfloor\log n\rfloor$ variables, the formula $\varphi_{n,K}$ is an
$O(K\log^{2}n)$-CNF of size $n^{O(K\log n)}$.
The formula $\varphi_{n,K}$ has resolution refutation of size $n^{O(K\log n)}$ and of width $O(K\log^{2}n)$
but every $\mathrm{Res}(\oplus)$ refutation of $\varphi_{n,K}$ has either size at least $2^{\lfloor n/40\log n\rfloor}$ or depth at least
$\Omega(K^{2}n\log n)$.\par



\par\medskip\noindent\textbf{Proof.}\par


Let $\Psi_{n}$ be a formula from Corollary 15.\par



Let us define $\varphi_{n,K}:=\Psi_{n}\circ\oplus_{K\lfloor\log n\rfloor}\circ Maj_{5}$.
Then $\varphi_{n,K}$ that contains $5Kn\lfloor\log n\rfloor$ variables, of the formula
$\varphi_{n,K}$ is in $O(K\log^{2}n)$-CNF and of size $n^{O(K\log n)}$.
By Lemma 6, the formula $\varphi_{n,K}$ has resolution refutation of size $n^{O(K\log n)}$ and of width $O(K\log^{2}n)$.\par



Let $w=\lfloor n/40\log n\rfloor$ and $h=\lfloor n^{2}/400\log^{2}n\rfloor$.
The formula $\Psi_{n}$ does not have resolution refutations of width at most $w$ and of depth at most $h$.
Let $s=K\lfloor\log n\rfloor$. It is easy to see that $s\geq 2$.\par



If $w<1$, then the conclusion of the theorem is trivial. Assume that $w\geq 1$.
Notice that $h+w\geq h+1\geq\frac{n^{2}}{400\log^{2}n}=\frac{n}{10\log^{3}n}\frac{n}{40\log
n}\cdot\log^{2}n\geq K^{2}\log^{2}nw\geq ws^{2}$, by Theorem 25, any $\mathrm{Res}(\oplus)$ refutation of $\varphi_{n,K}$ has either
size at least $2^{w}=2^{\lfloor n/40\log n\rfloor}$ or has depth at least $\frac{w^{2}s}{20}=\Omega(K^{2}n\log n)$.
$\hfill\blacktriangleleft$\par



\par\medskip\noindent\textbf{
Corollary 28.
}\par



For every $\delta>0$,
Depth-$\frac{n^{4/3}}{\log^{4/3+\delta}n}$ $\mathrm{Res}(\oplus)$ does not p-simulate resolution.\par



\par\medskip\noindent\textbf{Proof.}\par


Consider $K=\left\lfloor\sqrt{\frac{n}{10\log^{3}n}}\right\rfloor$ and the formula $\varphi_{n,K}$ from Theorem 27.
$\varphi_{n,K}$ contains
$m=\Theta\left(\frac{n^{3/2}}{\log^{1/2}n}\right)$ variables and
has resolution refutation of size at most $n^{O(\sqrt{n/\log n}})$ and the size of $\varphi_{n,K}$ is also $n^{O(\sqrt{n/\log n}})$ .
There is a constant $c$ such that every $\mathrm{Res}(\oplus)$ refutation of $\varphi_{n,K}$ of depth at most $cn^{2}/\log^{2}n$ has size at least $2^{\Omega\left(n/\log n\right)}$.\par



Notice that for $n$ large enough,
$\frac{m^{4/3}}{\log^{4/3+\delta}m}=\Theta\left(\frac{n^{2}}{\log^{2+\delta}n}\right)<cn^{2}/\log^{2}n$. Thus, for $n$ large enough every $\mathrm{Res}(\oplus)$ refutation of $\varphi_{n,K}$ of depth at most $\frac{m^{4/3}}{\log^{4/3+\delta}m}$ has size at least $2^{\Omega(n/\log n)}$. And $2^{\Omega(n/\log n)}$ can not be bounded by a polynomial in $n^{O\left(\sqrt{n/\log n}\right)}$.
$\hfill\blacktriangleleft$\par



\par\medskip\noindent\textbf{
Theorem 29.
}\par



Let $\psi_{n}$ be a family of unsatisfiable $O(1)$-CNF formulas such that $\psi_{n}$ has $n$ variables and the resolution width
of $\psi_{n}$ is $w(n)$.
For every natural $K\geq 2$ consider a formula $\Psi_{n,K}:=\psi_{n}\circ\oplus_{K}\circ Maj_{5}$;
it has $5nK$ variables, $\Psi_{n,K}$ is an $O(K)$-CNF formula of size at most $\mathrm{poly}(n)2^{K}$ and any $\mathrm{Res}(\oplus)$
refutations of $\Psi_{n}$ has either depth at least $\Omega(w(n)K^{2})$ or size at least $2^{\Omega(w(n))}$.\par



\par\medskip\noindent\textbf{Proof.}\par


Let $w(n)$ be the resolution width of $\psi_{n}$, take $s=K$. Let $h=w(n)s^{2}$.
There are no resolution refutations of $\psi_{n}$ of width at most $w(n)-1$ and depth $h$. $Maj_{5}$ is a 2-stifling gadget.
Then by Theorem 25, any $\mathrm{Res}(\oplus)$ refutation of $\Psi_{n}$ has either size at least $2^{w(n)-1}=2^{\Omega(n)}$
or depth at least $\frac{ws^{2}}{4\ell}=\Omega(nK^{2})$.
$\hfill\blacktriangleleft$\par


\par\medskip\noindent\textbf{
7 Open questions}\par


Two main avenues for improving our results are:\par


\begin{itemize}


\item 
1. 


Construct a polynomial-sized CNF formula that admits a polynomial-sized resolution
refutation, yet any $\mathrm{Res}(\oplus)$ refutation of it must have either superlinear depth or exponential size.
One approach to achieving this is by strengthening the supercritical tradeoff
between width and depth in resolution. Specifically, it suffices to construct an $O(1)$-CNF
formula with $n$ variables that has a resolution refutation of constant width,
but for which any resolution refutation must have either superlinear depth or width $\Omega(n)$.\par





\item 
2. 


Establish a truly supercritical tradeoff between size and depth for $\mathrm{Res}(\oplus)$,
in which the depth is superlinear with respect to the size of the formula.\par





\end{itemize}


\par\medskip\noindent\textbf{References}\par

\begin{itemize}


\item 
[1]

Miklós Ajtai.


The complexity of the pigeonhole principle.


Combinatorica, 14(4):417–433, 1994.


doi:10.1007/BF01302964.





\item 
[2]

Yaroslav Alekseev and Dmitry Itsykson.


Lifting to bounded-depth and regular resolutions over parities via
games.


In Michal Koucký and Nikhil Bansal, editors, Proceedings
of the 57th Annual ACM Symposium on Theory of Computing, STOC 2025,
Prague, Czechia, June 23-27, 2025, pages 584–595. ACM, 2025.


doi:10.1145/3717823.3718150.





\item 
[3]

Albert Atserias and Víctor Dalmau.


A combinatorial characterization of resolution width.


Journal of Computer and System Sciences, 74(3):323–334, 2008.


Computational Complexity 2003.


doi:10.1016/j.jcss.2007.06.025.





\item 
[4]

Paul Beame, Chris Beck, and Russell Impagliazzo.


Time-space trade-offs in resolution: Superpolynomial lower bounds for
superlinear space.


SIAM J. Comput., 45(4):1612–1645, 2016.


doi:10.1137/130914085.





\item 
[5]

Paul Beame and Sajin Koroth.


On Disperser/Lifting Properties of the Index and Inner-Product
Functions.


In Yael Tauman Kalai, editor, 14th Innovations in Theoretical
Computer Science Conference (ITCS 2023), volume 251 of Leibniz
International Proceedings in Informatics (LIPIcs), pages 14:1–14:17,
Dagstuhl, Germany, 2023. Schloss Dagstuhl – Leibniz-Zentrum für
Informatik.


doi:10.4230/LIPIcs.ITCS.2023.14.





\item 
[6]

Chris Beck, Jakob Nordström, and Bangsheng Tang.


Some trade-off results for polynomial calculus: extended abstract.


In Dan Boneh, Tim Roughgarden, and Joan Feigenbaum, editors, Symposium on Theory of Computing Conference, STOC’13, Palo Alto, CA, USA,
June 1-4, 2013, pages 813–822. ACM, 2013.


doi:10.1145/2488608.2488711.





\item 
[7]

Christoph Berkholz.


On the complexity of finding narrow proofs.


In 53rd Annual IEEE Symposium on Foundations of Computer
Science, FOCS 2012, New Brunswick, NJ, USA, October 20-23, 2012, pages
351–360. IEEE Computer Society, 2012.


doi:10.1109/FOCS.2012.48.





\item 
[8]

Christoph Berkholz and Jakob Nordström.


Supercritical space-width trade-offs for resolution.


SIAM J. Comput., 49(1):98–118, 2020.


doi:10.1137/16M1109072.





\item 
[9]

Sreejata Kishor Bhattacharya and Arkadev Chattopadhyay.


Exponential lower bounds on the size of reslin proofs of nearly
quadratic depth.


Electron. Colloquium Comput. Complex., TR25-106, 2025.


URL: https://eccc.weizmann.ac.il/report/2025/106.





\item 
[10]

Sreejata Kishor Bhattacharya, Arkadev Chattopadhyay, and Pavel Dvorák.


Exponential separation between powers of regular and general
resolution over parities.


In Rahul Santhanam, editor, 39th Computational Complexity
Conference, CCC 2024, July 22-25, 2024, Ann Arbor, MI, USA, volume 300
of LIPIcs, pages 23:1–23:32. Schloss Dagstuhl – Leibniz-Zentrum
für Informatik, 2024.


doi:10.4230/LIPIcs.CCC.2024.23.





\item 
[11]

Sam Buss and Neil Thapen.


A simple supercritical tradeoff between size and height in
resolution.


Information Processing Letters, 191:106589, 2026.


The preprint is available at
https://eccc.weizmann.ac.il/report/2024/001.


doi:10.1016/j.ipl.2025.106589.





\item 
[12]

Farzan Byramji and Russell Impagliazzo.


Lifting to randomized parity decision trees.


Electron. Colloquium Comput. Complex., TR24-202, 2024.


URL: https://eccc.weizmann.ac.il/report/2024/202.





\item 
[13]

Farzan Byramji and Russell Impagliazzo.


Lower bounds for the bit pigeonhole principle in bounded-depth
resolution over parities.


Electron. Colloquium Comput. Complex., TR25-118, 2025.


URL: https://eccc.weizmann.ac.il/report/2025/118.





\item 
[14]

Arkadev Chattopadhyay and Pavel Dvorák.


Super-critical trade-offs in resolution over parities via lifting.


In Srikanth Srinivasan, editor, 40th Computational Complexity
Conference, CCC 2025, August 5-8, 2025, Toronto, Canada, volume 339 of
LIPIcs, pages 24:1–24:19. Schloss Dagstuhl – Leibniz-Zentrum für
Informatik, 2025.


doi:10.4230/LIPIcs.CCC.2025.24.





\item 
[15]

Arkadev Chattopadhyay, Nikhil S. Mande, Swagato Sanyal, and Suhail Sherif.


Lifting to parity decision trees via stifling.


In Yael Tauman Kalai, editor, 14th Innovations in Theoretical
Computer Science Conference, ITCS 2023, January 10-13, 2023, MIT,
Cambridge, Massachusetts, USA, volume 251 of LIPIcs, pages
33:1–33:20. Schloss Dagstuhl – Leibniz-Zentrum für Informatik, 2023.


doi:10.4230/LIPIcs.ITCS.2023.33.





\item 
[16]

Susanna F. de Rezende, Noah Fleming, Duri Andrea Janett, Jakob Nordström,
and Shuo Pang.


Truly supercritical trade-offs for resolution, cutting planes,
monotone circuits, and weisfeiler–leman.


In Proceedings of the 57th Annual ACM Symposium on Theory of
Computing, STOC ’25, pages 1371–1382, New York, NY, USA, 2025. Association
for Computing Machinery.


doi:10.1145/3717823.3718271.





\item 
[17]

Klim Efremenko, Michal Garlík, and Dmitry Itsykson.


Lower bounds for regular resolution over parities.


SIAM J. Comput., 54(4):887–915, 2025.


Preliminary version appeared in Proceedings of STOC 2024.


doi:10.1137/24M1696640.





\item 
[18]

Klim Efremenko and Dmitry Itsykson.


Amortized closure and its applications in lifting for resolution over
parities.


In Srikanth Srinivasan, editor, 40th Computational Complexity
Conference, CCC 2025, August 5-8, 2025, Toronto, Canada, volume 339 of
LIPIcs, pages 8:1–8:24. Schloss Dagstuhl – Leibniz-Zentrum für
Informatik, 2025.


doi:10.4230/LIPIcs.CCC.2025.8.





\item 
[19]

Noah Fleming, Toniann Pitassi, and Robert Robere.


Extremely Deep Proofs.


In Mark Braverman, editor, 13th Innovations in Theoretical
Computer Science Conference (ITCS 2022), volume 215 of Leibniz
International Proceedings in Informatics (LIPIcs), pages 70:1–70:23,
Dagstuhl, Germany, 2022. Schloss Dagstuhl – Leibniz-Zentrum für
Informatik.


doi:10.4230/LIPIcs.ITCS.2022.70.





\item 
[20]

Michal Garlík and Leszek Aleksander Kolodziejczyk.


Some subsystems of constant-depth frege with parity.


ACM Trans. Comput. Log., 19(4):29:1–29:34, 2018.


doi:10.1145/3243126.





\item 
[21]

Mika Göös, Gilbert Maystre, Kilian Risse, and Dmitry Sokolov.


Supercritical tradeoffs for monotone circuits.


In Michal Koucký and Nikhil Bansal, editors, Proceedings
of the 57th Annual ACM Symposium on Theory of Computing, STOC 2025,
Prague, Czechia, June 23-27, 2025, pages 1359–1370. ACM, 2025.


doi:10.1145/3717823.3718229.





\item 
[22]

Svyatoslav Gryaznov, Sergei Ovcharov, and Artur Riazanov.


Resolution over linear equations: Combinatorial games for tree-like
size and space.


ACM Trans. Comput. Theory, 16(3):15:1–15:15, 2024.


doi:10.1145/3675415.





\item 
[23]

Dmitry Itsykson and Artur Riazanov.


Proof complexity of natural formulas via communication arguments.


In Valentine Kabanets, editor, 36th Computational Complexity
Conference, CCC 2021, July 20-23, 2021, Toronto, Ontario, Canada (Virtual
Conference), volume 200 of LIPIcs, pages 3:1–3:34. Schloss Dagstuhl
– Leibniz-Zentrum für Informatik, 2021.


doi:10.4230/LIPIcs.CCC.2021.3.





\item 
[24]

Dmitry Itsykson and Dmitry Sokolov.


Lower bounds for splittings by linear combinations.


In Erzsébet Csuhaj-Varjú, Martin Dietzfelbinger, and
Zoltán Ésik, editors, Mathematical Foundations of Computer
Science 2014 - 39th International Symposium, MFCS 2014, Budapest, Hungary,
August 25-29, 2014. Proceedings, Part II, volume 8635 of Lecture
Notes in Computer Science, pages 372–383. Springer, 2014.


doi:10.1007/978-3-662-44465-8\_32.





\item 
[25]

Dmitry Itsykson and Dmitry Sokolov.


Resolution over linear equations modulo two.


Ann. Pure Appl. Log., 171(1), 2020.


doi:10.1016/J.APAL.2019.102722.





\item 
[26]

Jan Krajíček.


Randomized feasible interpolation and monotone circuits with a local
oracle.


J. Math. Log., 18(2):1850012:1–1850012:27, 2018.


doi:10.1142/S0219061318500125.





\item 
[27]

Pavel Pudlák and Russell Impagliazzo.


A lower bound for DLL algorithms for k-sat (preliminary
version).


In David B. Shmoys, editor, Proceedings of the Eleventh Annual
ACM-SIAM Symposium on Discrete Algorithms, January 9-11, 2000, San
Francisco, CA, USA, pages 128–136. ACM/SIAM, 2000.


URL: http://dl.acm.org/citation.cfm?id=338219.338244.





\item 
[28]

A. A. Razborov.


Lower bounds on the size of bounded depth circuits over a complete
basis with logical addition.


Mat. Zametki, 41:598–607, 1987.





\item 
[29]

Alexander A. Razborov.


A new kind of tradeoffs in propositional proof complexity.


J. ACM, 63(2):16:1–16:14, 2016.


doi:10.1145/2858790.





\item 
[30]

Alexander A. Razborov.


On the width of semialgebraic proofs and algorithms.


Math. Oper. Res., 42(4):1106–1134, 2017.


doi:10.1287/MOOR.2016.0840.





\item 
[31]

Alexander A. Razborov.


On space and depth in resolution.


Comput. Complex., 27(3):511–559, 2018.


doi:10.1007/S00037-017-0163-1.





\item 
[32]

Roman Smolensky.


Algebraic methods in the theory of lower bounds for boolean circuit
complexity.


In Alfred V. Aho, editor, Proceedings of the 19th Annual ACM
Symposium on Theory of Computing, 1987, New York, New York, USA, pages
77–82. ACM, 1987.


doi:10.1145/28395.28404.





\item 
[33]

Alasdair Urquhart.


Hard examples for resolution.


Journal of the ACM, 34(1):209–219, 1987.


doi:10.1145/7531.8928.





\item 
[34]

Alasdair Urquhart.


The depth of resolution proofs.


Stud Logica, 99(1-3):349–364, 2011.


doi:10.1007/S11225-011-9356-9.





\end{itemize}
